
% ════════════════════════════════════════════
% 模型建立（5.2.3.模型建立.tex，三级子块）—— 写作提示词见本文件末尾【本章节写作提示词】。
% ════════════════════════════════════════════
\subsubsection{模型建立}

针对问题二，本节将建立两站交会定位精度模型与并发告警车辆调度模型。具体过程如下。

\textbf{步骤1：交会定位误差传播模型}

设两观测站 $S_1$、$S_2$ 到目标水平投影的距离分别为 $r_1$、$r_2$，方位角测量误差 $\varepsilon_1,\varepsilon_2$ 相互独立且服从 $N(0,\sigma^2)$。在目标处建立局部坐标，方位误差 $\varepsilon_i$ 使第 $i$ 条方位线发生横向平移 $r_i\varepsilon_i$，两条平移后的方位线交点为定位结果。设两方位线交会角为 $\gamma$，经线性化推导可得位置误差向量沿两方位线法向的分量满足：
\begin{equation}
    \bm{d}\cdot \bm{v}_1 = r_1\varepsilon_1, \qquad \bm{d}\cdot \bm{v}_2 = r_2\varepsilon_2,
\end{equation}
其中，$\bm{v}_i$ 为第 $i$ 条方位线的单位法向量，且 $\bm{v}_1\cdot\bm{v}_2=\cos\gamma$。解得位置误差为：
\begin{equation}
    \bm{d}=\frac{r_2\varepsilon_2-r_1\varepsilon_1\cos\gamma}{\sin\gamma}\,\bm{u}_1 + r_1\varepsilon_1\,\bm{v}_1 .
\end{equation}
对模长取期望，并利用 $E[\varepsilon_i^2]=\sigma^2$、$E[\varepsilon_1\varepsilon_2]=0$，得定位误差均方根为：
\begin{equation}
    \sigma_p=\sqrt{E[\lVert\bm{d}\rVert^2]}=\sigma\,\frac{\sqrt{r_1^2+r_2^2}}{\lvert\sin\gamma\rvert} .
\end{equation}
该式表明：距离越近、交会角越接近 90 度，定位越精确；交会角趋于 0 度或 180 度时 $\lvert\sin\gamma\rvert\to 0$，误差发散。这与多站测向交叉定位精度分析的一般结论一致\cite{jiang2011shuxue,wang2007cexiang,liu2008cexiang}。

\textbf{步骤2：精度指标与退化讨论}

由式(12)定义交会定位精度指标为：
\begin{equation}
    P=\frac{\lvert\sin\gamma\rvert}{\sqrt{r_1^2+r_2^2}},
\end{equation}
$P$ 越大定位越精确。当 $\gamma\to 0^\circ$ 或 $180^\circ$ 时，两方位线近乎平行，微小误差即可导致交点大幅偏移，模型退化；当站点位于可观测区域边界时，俯仰角处于 80 度临界或视线恰好擦过建筑，轻微扰动即可能破坏可观测性。处理策略为：对交会角施加下限 $\gamma\in[40^\circ,140^\circ]$，并让站点与边界保持净距余量。

取 $\gamma_{\min}=40^\circ$ 的依据是：误差放大因子 $1/\lvert\sin\gamma\rvert$ 在 $40^\circ$ 处为 1.56，即最坏情况下交会精度损失不超过 56\%；放宽至 $20^\circ$ 时该因子升至 2.92，收紧至 $50^\circ$ 则收益甚微。数值检验表明，$\gamma\ge40^\circ$ 时三个告警的 $\sigma_p$ 均不超过 162 cm，最大观测开始时间为 12.063 s；若放宽至 $20^\circ$，最大观测开始时间降至 10.830 s，但 $\sigma_p$ 恶化至 410 cm，且最近指派启发式与最优解重合，调度问题退化为“原地观测”。若将该角度约束替换为直接精度约束 $\sigma_p\le2$ m，所得最优解与 12.063 s 完全一致，说明二者在本算例中等价；进一步收紧至 1.5 m 与 1.2 m 时，最大观测开始时间分别升至 15.846 s 与 22.509 s。可见 2 m（对应 $40^\circ$）位于时间-精度权衡的拐点附近，是兼顾效率与精度的合理取值。

\textbf{步骤3：最优站对优化模型}

在车辆行驶时间不超过 30 s 的约束下，选择 2 台车并确定 T2 的观测站对，使精度指标最大：
\begin{equation}
    \max_{S_1,S_2,v_a,v_b} \; P \quad \text{s.t.} \quad
    \begin{cases}
    S_1,S_2 \text{ 满足三项可观测约束},\\
    \dfrac{\lVert S_1-V_a\rVert}{v_s}\le 30,\; \dfrac{\lVert S_2-V_b\rVert}{v_s}\le 30,\\
    v_a\neq v_b .
    \end{cases}
\end{equation}
其中，$v_s=3$ m/s 为行驶速度，$V_a$、$V_b$ 为两车初始位置。

下面证明该最优站对无需连续搜索即可全局最优。设俯仰约束给出的最小水平距离为 $r_{\min}$，即两站均满足 $r_1,r_2\ge r_{\min}$；又 $\lvert\sin\gamma\rvert\le 1$ 恒成立。由误差传播公式 $\sigma_p=\sigma\,\sqrt{r_1^2+r_2^2}/\lvert\sin\gamma\rvert$ 可得下界：
\begin{equation}
    \sigma_p=\sigma\,\frac{\sqrt{r_1^2+r_2^2}}{\lvert\sin\gamma\rvert}\ge \sigma\sqrt{2}\,r_{\min},
\end{equation}
等号当且仅当 $r_1=r_2=r_{\min}$ 且 $\lvert\sin\gamma\rvert=1$（即 $\gamma=90^\circ$）时成立，对应“正西 + 正北、距离取 $r_{\min}$”的构型。因此，只需校验该构型是否满足可观测与行驶时间约束即可判定全局最优，无需对站对做连续搜索。

\textbf{步骤4：并发告警车辆调度模型}

对并发告警 $i\in\{1,2,3\}$，设车辆 $v$ 沿任务序列行驶，其到第 $i$ 告警观测站的到站时刻与云台转动时长分别为：
\begin{equation}
    t_{\text{arr}}(v,i)=t_{\text{free}}(v)+\frac{d(v,S_{v,i})}{v_s}, \qquad
    t_{\text{rot}}(v,i)=\frac{\max(\Delta\alpha_{v,i},\,\lvert\varphi_{v,i}\rvert)}{\omega},
\end{equation}
其中，$d$ 为行驶距离，$\omega=30$ 度/s 为云台角速度。车辆对告警 $i$ 的就绪时刻为：
\begin{equation}
    t_{\text{ready}}(v,i)=t_{\text{arr}}(v,i)+t_{\text{rot}}(v,i).
\end{equation}
告警 $i$ 的观测开始时刻为两车就绪时刻的较大者：
\begin{equation}
    t_{\text{start}}(i)=\max\big(t_{\text{ready}}(v_a,i),\,t_{\text{ready}}(v_b,i)\big).
\end{equation}
目标是最小化最大观测开始时间：
\begin{equation}
    \min \; \max_i t_{\text{start}}(i).
\end{equation}
每辆车完成观测后需以 $t_{\text{rot}}$ 复位云台，再前往下一观测站，其时间轴据此递推。

综上所述，本节完成了交会定位精度模型与并发调度模型的建立，接下来将基于此进行求解。
